% Paket und Style fuer Code-Listings
\usepackage{listings}
\usepackage{color}
\renewcommand{\lstlistingname}{Listing}

\definecolor{codebg}{RGB}{245,245,245}
\definecolor{codekeyword}{RGB}{0,0,150}
\definecolor{codecomment}{RGB}{100,100,100}
\definecolor{codestring}{RGB}{163,21,21}
\definecolor{codetorch}{RGB}{170,60,0}

% PyTorch/NumPy-Bezeichner, die im Quellcode der Kapitel als eigene
% Keyword-Klasse (Klassen, Module, Tensor-Methoden) hervorgehoben werden
\lstdefinestyle{bachelorarbeit}{
    backgroundcolor=\color{codebg},
    basicstyle=\ttfamily\small,
    keywordstyle=\color{codekeyword}\bfseries,
    keywordstyle=[2]\color{codetorch}\bfseries,
    commentstyle=\color{codecomment}\itshape,
    stringstyle=\color{codestring},
    numbers=left,
    numberstyle=\tiny\color{gray},
    numbersep=8pt,
    frame=single,
    breaklines=true,
    breakatwhitespace=false,
    showstringspaces=false,
    tabsize=4,
    captionpos=b,
    morekeywords=[2]{
        torch, nn, np,
        DataLoader, Dataset, GRURegressor,
        Adam, MSELoss,
        cuda, is_available, manual_seed, seed,
        no_grad, clip_grad_norm_, utils,
        state_dict, load_state_dict, parameters,
        zero_grad, backward, step, train, eval,
        to, from_numpy, tensor, detach, cpu, clone,
        float32
    },
    literate=
        {ä}{{\"a}}1 {ö}{{\"o}}1 {ü}{{\"u}}1
        {Ä}{{\"A}}1 {Ö}{{\"O}}1 {Ü}{{\"U}}1
        {ß}{{\ss}}1,
}

\lstset{style=bachelorarbeit}
